%% Jaren Untuña --- Curriculum Vitae (English)
%% Compile:  tectonic cv-en.tex
%% Plantilla propia: una columna, sin iconos, sin barras de "nivel de skill".
%% Está pensada para pasar un ATS y para que una persona lo lea en 20 segundos.

\documentclass[10.5pt,a4paper]{article}

\usepackage[T1]{fontenc}
\usepackage{tgheros}
\renewcommand*\familydefault{\sfdefault}
\usepackage[utf8]{inputenc}
\usepackage[english]{babel}
\usepackage[
  a4paper,
  top=1.5cm, bottom=1.5cm, left=1.6cm, right=1.6cm,
  footskip=0.8cm
]{geometry}
\usepackage{titlesec}
\usepackage{enumitem}
\usepackage{xcolor}
\usepackage{hyperref}
\usepackage{tabularx}
\usepackage{microtype}

\definecolor{ink}{HTML}{1A1A1A}     % casi negro, para nombre y titulares
\definecolor{muted}{HTML}{5F6670}   % gris medio, texto secundario
\definecolor{soft}{HTML}{8A8F98}    % gris claro, etiquetas y fechas
\definecolor{rule}{HTML}{D8DBDF}    % gris muy claro, lineas divisorias
\definecolor{link}{HTML}{2B2B2B}    % gris oscuro para enlaces (sin azul)

\hypersetup{
  colorlinks=true,
  urlcolor=link,
  linkcolor=link,
  pdftitle={Johan Untuna - Software Engineer - CV},
  pdfauthor={Johan Untuna},
  pdfsubject={Curriculum Vitae},
  pdfkeywords={software engineer, python, django, data pipelines, api integration, security}
}

\color{ink}
\pagestyle{empty}
\setlength{\parindent}{0pt}
\setlength{\emergencystretch}{2em}

%% ---------------------------------------------------------------- headings
\titleformat{\section}
  {\normalfont\large\bfseries\scshape}
  {}{0pt}{}
  [\vspace{-0.5em}{\color{rule}\rule{\linewidth}{0.4pt}}]
\titlespacing*{\section}{0pt}{1.15em}{0.6em}

%% Entry header: title on the left, dates flush right.
\newcommand{\entry}[2]{%
  \begin{tabularx}{\linewidth}{@{}Xr@{}}
    \textbf{#1} & {\small\color{soft}#2}
  \end{tabularx}%
  \vspace{0.15em}
}

%% One-line subtitle under an entry header.
\newcommand{\subentry}[1]{{\small\color{muted}#1}\vspace{0.25em}}

\newenvironment{points}
  {\begin{itemize}[leftmargin=1.1em, itemsep=0.18em, topsep=0.25em, parsep=0pt, label=\textcolor{soft}{\footnotesize\textbullet}]}
  {\end{itemize}}

%% Skill row: bold category, then the list.
\newcommand{\skillrow}[2]{%
  \begin{tabularx}{\linewidth}{@{}p{3.5cm}X@{}}
    \textbf{#1} & #2
  \end{tabularx}%
  \vspace{0.15em}
}

\begin{document}

%% ------------------------------------------------------------------ header
{\Huge\bfseries Johan Untuña}\\[0.35em]
{\large\color{muted} Software Engineer --- Data Pipelines, API Integration \& Applied AI}\\[0.7em]
{\small
Latacunga, Ecuador (GMT-5) \textbullet{} Open to remote \\
\href{mailto:jaren22uv@gmail.com}{jaren22uv@gmail.com} \textbullet{}
\href{https://wa.me/593963607760}{+593 96 360 7760} \\
\href{https://github.com/JohanUV}{github.com/JohanUV} \textbullet{}
\href{https://www.linkedin.com/in/jaren-untu\%C3\%B1a-valiente-762799211/}{LinkedIn} \textbullet{}
\href{https://tryhackme.com/p/jaren22uv}{TryHackMe}
}

\vspace{0.45em}
{\color{ink}\rule{\linewidth}{1.1pt}}

%% ----------------------------------------------------------------- summary
\section*{Profile}
Software engineering student (5th semester, ESPE Latacunga) who builds and ships complete
systems rather than exercises. My work centres on getting data out of sources that resist
being read --- undocumented government APIs, portals behind captchas, providers with
incompatible schemas --- and making the result trustworthy enough to decide on. I design with
access control, audit trails and data-protection rules in from the start. Two of my projects
were commissioned by local businesses; the rest are public on GitHub with the architecture
documented. Moving toward offensive security and OSINT.

%% ---------------------------------------------------------------- projects
\section*{Selected Projects}

\entry{Verum --- Pre-hire Due Diligence Platform}{2026}
\subentry{Python, Flask, React, REST, RBAC \textbullet{} \href{https://github.com/JohanUV/Verum-EC}{github.com/JohanUV/Verum-EC}}
\begin{points}
  \item Integrated \textbf{nine Ecuadorian public-record sources} --- judiciary, tax authority,
        child-support registry, traffic authority and OFAC/UN sanctions lists --- into a single
        consolidated risk report queried by national ID.
  \item Solved captcha-protected JSF sources with a \textbf{human relay} rather than defeating
        them: the app surfaces the captcha to the user and continues the authenticated session,
        keeping the integration within what each source permits.
  \item Labelled every source with an explicit reliability status (live API, captcha relay,
        assisted) so the report never overstates its own coverage.
  \item Built for Ecuador's data-protection law: role-based permissions, full audit history of
        who queried whom, and PDF export stamped with folio and watermark for traceability.
\end{points}

\entry{Job Radar --- Automated Job-Intelligence Pipeline}{2026}
\subentry{Django, DRF, PostgreSQL, n8n, React, Docker, Gemini \textbullet{} \href{https://github.com/JohanUV/Job-Radar}{github.com/JohanUV/Job-Radar}}
\begin{points}
  \item Scheduled pipeline collecting vacancies from three job APIs every six hours,
        normalising each provider into a shared schema through per-source connectors.
  \item Made duplicates \textbf{structurally impossible} by enforcing a unique SHA-256 hash of
        the vacancy URL at the database level, so re-running the pipeline cannot corrupt data
        regardless of concurrency.
  \item Split orchestration (n8n: scheduling, parallel fetch, retries) from business rules
        (Django: validation, deduplication, persistence) --- adding a source touches no backend code.
  \item Added an LLM stage scoring CV-to-vacancy affinity 0--100 with Gemini, made pull-based
        and idempotent so it can crash and resume without double-billing the model; stores the
        reasoning and model version alongside every score.
\end{points}

\entry{Judicial Records Analysis --- Reverse-Engineered Government API}{2026}
\subentry{Python, HTTP traffic analysis, API specification \textbullet{} \href{https://github.com/JohanUV/Sistema-Analisis-Judicial}{github.com/JohanUV/Sistema-Analisis-Judicial}}
\begin{points}
  \item Reconstructed the undocumented service contract behind Ecuador's judicial records
        portal from observed traffic: two endpoints, full payload shape, pagination and the
        headers the server actually validates.
  \item Identified that the service validates \textbf{provenance} (Origin/Referer) rather than
        identity, and that its captcha field is accepted as a constant --- a client-side-trust
        finding documented at the service layer.
  \item Wrote the result up as a reusable API reference instead of leaving it implicit in
        scraper code.
\end{points}

\entry{Class Attention --- Classroom Attention Detection}{2026}
\subentry{Python, machine learning, computer vision, privacy by design}
\begin{points}
  \item ML system estimating student attention during a class from video, validated on real
        classroom sessions, giving teachers a signal they can act on mid-lesson.
  \item Designed so the useful output is an \textbf{aggregate over the room}, never a per-student
        record --- removing the privacy risk at the architecture level rather than by policy.
  \item Identity protection treated as a constraint from the first version, since the subjects
        are minors and a leak would be unrecoverable.
\end{points}

%% ------------------------------------------------------------ client work
\section*{Client Work}

\entry{Cottullari --- Tourism Operator, Latacunga}{2026}
\subentry{React, React Router, Vite, PHP, MySQL \textbullet{} \href{https://github.com/JohanUV/cottullari-react}{github.com/JohanUV/cottullari-react}}
\begin{points}
  \item Built a seven-route single-page application for a tourism company: fleet, services,
        destinations and quote requests, with a form persisting enquiries to MySQL.
  \item Resolved all five destination pages through one parameterised route, so adding a
        destination is data rather than a new component.
\end{points}

\entry{Mashka Box --- CrossFit Gym, Latacunga}{2026}
\subentry{HTML, CSS, responsive \textbullet{} \href{https://github.com/JohanUV/mashka-box}{github.com/JohanUV/mashka-box}}
\begin{points}
  \item Mobile-first landing page built to convert a visit into a trial class; hand-written
        HTML and CSS with no framework, because a single page with no application state should
        not make the visitor download one.
\end{points}

%% ------------------------------------------------------------------ skills
\section*{Technical Skills}
\skillrow{Languages}{Python, TypeScript, JavaScript, SQL, HTML, CSS}
\skillrow{Backend}{Django, Django REST Framework, Flask, REST API design}
\skillrow{Frontend}{React, React Router, Vite, Astro, Tailwind CSS}
\skillrow{Data}{PostgreSQL, SQLite, Drizzle ORM, Pandas, web scraping, schema normalisation}
\skillrow{AI \& ML}{Gemini API, OpenAI API, Ollama, OpenCV, scikit-learn, prompt design}
\skillrow{Automation}{n8n, scheduled pipelines, connector architectures, Telegram bots}
\skillrow{Infrastructure}{Linux, Bash, Docker, Docker Compose, Terraform, LocalStack, AWS Lambda, Vercel, Git}
\skillrow{Security}{Access control (RBAC), audit trails, API-key authentication, data-protection compliance}
\skillrow{Learning}{Offensive security, Burp Suite, Nmap, network protocols, OSINT}

%% --------------------------------------------------------------- education
\section*{Education}
\entry{Universidad de las Fuerzas Armadas ESPE --- Latacunga}{2024 --- 2028 (expected)}
\subentry{B.Eng. Software Engineering \textbullet{} Currently 5th semester}

%% --------------------------------------------------------------- languages
\section*{Languages}
\skillrow{Spanish}{Native}
\skillrow{English}{B2 --- professional working proficiency, comfortable in technical discussion}

\end{document}
